%
% Der Funktionekörper
%
\subsection{Der Funktionenkörper $\mathbb{Q}(x,\vartheta)$
\label{buch:intalg:wurzel:subsection:koerper}}
Der Körper $\mathbb{Q}(z,\vartheta)$ besteht aus den Funktionen
\[
f(x)
=
\frac{
\tilde{p}_0 + \tilde{p}_1\vartheta + \tilde{p}_2\vartheta^2
+ \ldots + \tilde{p}_n\vartheta^n
}{
\tilde{q}_0 + \tilde{q}_1\vartheta + \tilde{q}_2\vartheta^2
+ \ldots + \tilde{q}_m\vartheta^m
}
\]
mit $p_k\in\mathbb{Q}(x)$, $k=1,\dots,n$, und $q_k\in\mathbb{Q}(x)$,
$k=1,\dots,m$.
Da $\vartheta$ eine algebraisches Element über $\mathbb{Q}$ ist, welches
die Identität
\[
\vartheta^2 = ax^2 +bx +c
\]
mit $a,b,c\in\mathbb{Q}$ erfüllt, können die höheren Potenzen
$\vartheta^{2l}$ und $\vartheta^{2l+1}$ durch Potenzen des Radikanden
ersetze werden:
\begin{align*}
\vartheta^{2l}
&=
(ax^2 + bx + c)^l
&&\text{und}&
\vartheta^{2l+1}
&=
(ax^2 + bx + c)^l \vartheta.
\end{align*}
Es ist daher möglich, die Funktion $f(x)$ in der einfacheren Form
\begin{equation}
f(x)
=
\frac{p_0+p_1\vartheta}{q_0+q_1\vartheta}
\label{buch:algint:wurzel:eqn:funktion}
\end{equation}
zu schreiben.
Gesucht ist also eine Stammfunktion von
\[
f
\in
\mathbb{Q}(x,\vartheta)
=
\biggl\{
\frac{p_0+p_1\vartheta}{q_0+q_1\vartheta}
\;
\bigg|
\;
p_k,q_k\in\mathbb{Q}(x), k=0,1
\biggr\}
\]
als elementare Funktion.
Da schon die Stammfunktionen von rationalen Funktionen nicht unbedingt
rational sein müssen, kann man nicht erwarten, dass eine Stammfunktion
einer Funktion in $\mathbb{Q}(x,\vartheta)$ in diesem Körper liegen wird.
Es werden also Erweiterungen des Körpers um weitere Funktionen nötig sein,
um das Problem lösen zu können.
